% Notes with margin notes, key boxes and a checklist.
\documentclass[10pt]{article}
\usepackage{tinfoil-base}
% ---- PERSONALIZATION --------------------------------------------------
\newcommand{\NotesTitle}{Distributed Systems, Lecture 7}
\newcommand{\NotesSubtitle}{Consensus and the FLP result}
\newcommand{\NotesMeta}{21 November 2026 \ $\bullet$\ Prof.\ Haddad \ $\bullet$\ Week 7}
% -----------------------------------------------------------------------
\geometry{left=2cm,right=5.4cm,marginparwidth=3.6cm,marginparsep=0.6cm}
\pagestyle{empty}
\setcounter{secnumdepth}{0}
\newcommand{\sidenote}[1]{\marginpar{\footnotesize\color{muted}\raggedright #1}}
\newcommand{\todo}[1]{\item[$\square$] #1}
\newcommand{\done}[1]{\item[$\blacksquare$] #1}
\begin{document}
\docheader{\NotesTitle}{\NotesSubtitle}
\vspace{-1em}{\small\color{muted}\NotesMeta}\par\vspace{0.5em}

% ---- BODY ----
\section{Setting}
A set of $n$ processes communicate by message passing. Messages may be delayed arbitrarily but are eventually delivered. Up to $f$ processes may crash.\sidenote{Asynchronous model: no bound on message delay or relative process speed.}

\begin{keybox}{Definition: consensus}
Every correct process decides a value such that (1) all decisions are equal, (2) the decided value was proposed by some process, and (3) every correct process eventually decides.
\end{keybox}

\section{The FLP impossibility}
In the asynchronous model, no deterministic protocol solves consensus if even one process may crash.\sidenote{Fischer, Lynch and Paterson, 1985.}

\subsection{Proof sketch}
\begin{itemize}
  \item Some initial configuration is \emph{bivalent}: both outcomes remain possible.
  \item From any bivalent configuration, there is a step that leads to another bivalent configuration.
  \item Hence an infinite non-deciding run exists.
\end{itemize}

\section{Ways around it}
\begin{tabular}{@{}ll@{}}\toprule
\lab{Approach} & \lab{What it gives up}\\\midrule
Partial synchrony (Paxos, Raft) & Liveness during asynchrony\\
Randomisation (Ben-Or) & Determinism\\
Failure detectors & Pure asynchrony\\\bottomrule
\end{tabular}

\section{To do}
\begin{itemize}
  \done{Read FLP, Sections 1 to 3}
  \todo{Work through the bivalence lemma by hand}
  \todo{Compare with the Raft paper's safety argument}
\end{itemize}
% ---- END BODY ----
\end{document}
